\begin{proposition}[Monodromy and the smooth locus of an irreducible fibre]\hypertarget{tag-02WK}{}\label{night:P4.13}\label{stmt:P4.13}
\StatementRecord{night:P4.13}{proved}{02WK}{the definitions and constructions stated in the proof}

Let a proper holomorphic degeneration of two-dimensional complex
tori have smooth total threefold and central fibre $mS$ with
irreducible reduced support. If $A$ is the monodromy on $H_1$ of
a nearby torus, then
\[
 m e(S_{\mathrm{reg}})=\det(I-A^m),\qquad
 0\le e(S_{\mathrm{reg}})\le16/m.
\]
\end{proposition}
\begin{proof}

\begin{enumerate}[label=\arabic*.,leftmargin=*]
\item Choose a proper holomorphic resolution of the pair
$\pi:U'\to U$ which is an isomorphism away from $S$ and over
$S_{\mathrm{reg}}$. The total central divisor has normal crossings.
The strict transform has multiplicity $m$, and each exceptional
component has multiplicity at least $2m$. Indeed, its first centre
lies over a singular point of the reduced divisor, whose order is
at least two; subsequent centres either retain such an exceptional
multiplicity or add further positive multiplicities.

\item At a point on $k$ branches,
with local equation $z_1^{n_1}\cdots z_k^{n_k}=t$, the nearby
fibre retracts onto a disjoint union of real $(k-1)$-tori.
Monodromy permutes these components and translates any component
that it preserves. For $k\ge2$ the alternating cohomological
trace is zero. On a single branch of multiplicity $n$ there are
$n$ contractible nearby components in a cyclic permutation;
the $d$th iterate has trace $n$ when $n\mid d$, and zero otherwise.
Let $F$ be the nearby torus and $T$ its induced cohomology monodromy. Alternating traces in the nearby-cycle spectral sequence, with a
finite stratification of the compact central divisor, give
\[
 L(T^d;H^*(F;\C))=\sum_{n_j\mid d}n_j e(E_j^\circ).
\]
This is the proper analytic trace formula of \citet[Theorem~1, p.~235; \S\S2--3, pp.~237--245]{ACampo1975}, obtained from the displayed local traces and the nearby-cycle spectral sequence.
For $d=m$ the strict transform is the sole contribution, and its
open stratum is $S_{\mathrm{reg}}$.

\item For a common multiple $N$
of the $n_j$, the $N$th power of monodromy acts trivially on all local
nearby-cohomology sheaves, hence on the $E_2$ page. This $N$th power
acts unipotently on the filtered abutment. Every
eigenvalue of $A$ is a root of unity. Under the evaluation pairing $H^1(F;\C)=H_1(F;\C)^\vee$, the induced cohomology action is $A^T$, since $(T\alpha)(c)=\alpha(Ac)$. The cohomology of the nearby torus $F$ is the exterior algebra on $H^1(F;\C)$, so its Lefschetz number is
$\det(I-(A^m)^T)=\det(I-A^m)$. Complex conjugate eigenvalues contribute
$|1-\lambda^m|^2$; real eigenvalues are $1$ or $-1$ and
contribute zero or two. Thus this determinant is nonnegative
and at most $2^4$.

\end{enumerate}
\end{proof}